\begin{lemma}[Transport of the retained beta choice to condition (9)]
Fix the graph, palette, edge, input, displayed tree, generating coloring,
ETT, base, counterexample \(C\), minimum context, and two replacement
presentations \(P,R\) relative to these same original data.
Let \(S\) be a \noderef{P2880_T1_SplitSwapData} for
\(P,\alpha,\beta,\sigma\), and assume \(R.\mathit{coloring}=\sigma\).
Use a subscript \(P\) or \(R\) for the following data of that
presentation: its last starred stage \(A\), last stored set \(D\),
last Gamma union \(\Gamma\), terminal ladder \(T_n\), and
\[
 M_0=\bigcup_{v\in V(\mathit{replacementBase.core})\setminus V(T_n)}
 \noderef{MissingColors}(G,\mathit{coloring},v),\quad
 U_j=\{c:\exists r\in\operatorname{Fin}|T|,+ |A|\le r<|A|+j,\ \mathit{coloring}(T.\mathit{edge}(r))=c\}.
\]
Here \(j=\noderef{HierarchyPathNumber}(C.\mathit{hierarchy})\) and
\(B=\noderef{HierarchyFullTailPrefix}(C,j-2)\).
Assume \(\beta\ne\alpha\), \(\beta\in\overline{P.\mathit{coloring}}(B)\),
\(\beta\notin\Gamma_P\cup U_{j,P}\), and either
\(\beta\notin M_{0,P}\cup D_P\), or
\(\beta\in D_P\) and
\((P.\mathit{counterexample.good}.\Gamma^q_\beta\cup\{\beta\})
\cap U_{j,P}=\varnothing\).
Then \(\beta\notin\Gamma_R\), and either
\[
 \beta\notin M_{0,R},\quad\beta\notin U_{j,R},\quad
 \beta\notin\Gamma_R,\quad\beta\notin D_R,
\]
or
\[
 \beta\in D_R\cap\overline{R.\mathit{coloring}}(B),\qquad
 (R.\mathit{counterexample.good}.\Gamma^q_\beta\cup\{\beta\})
 \cap U_{j,R}=\varnothing.
\]
\end{lemma}
\begin{proof}
By \noderef{HierarchyReplacementPresentation}, both \(P\) and \(R\)
have \noderef{SameGoodHierarchy} with the fixed original good
hierarchy. Expand that definition and its \noderef{SameHierarchyStages}
clause. Each identifies its base core, divider count, and entire
dependent stage family with the original ones. Their transitive
equalities give equality of the two cores and divider counts.
\noderef{P2752_T2_HierarchyMeasures}, applied separately to the two
original-relative same-stage witnesses, gives \(A_P=A_R\).
The stored-set equality clauses at the last index give \(D_P=D_R\);
the dependent cast of a last index is the last index because its value
is the equal divider count. For each \(\eta\in D_P\), compose the
two Gamma-pair equalities through the original pair. This gives equality
of the last Gamma pairs on the common stored set, and taking their
finite unions gives \(\Gamma_P=\Gamma_R\).

We also identify the actual terminal ladder trees, not just the two
displayed final trees. For each original-relative
\noderef{CorrespondingExtendedTashkinovTree} certificate, eliminate
its rung equality, then its series-presentation heterogeneous equality.
At the positive index \(i=n-1\), its ladder agreement identifies the
replacement ladder at \(i+1=n\) with the original terminal ladder.
The index exists by the positive-rung field of
\noderef{ExtendedTashkinovTree}. Thus both actual terminal ladders
are the same exact original tree, and \(T_{n,P}=T_{n,R}\).
Consequently the vertex difference defining the two \(M_0\)'s is
the same finite set. It lies inside the common base core, which is a
prefix of \(X\) by \(S\)'s core-in-split field.
\noderef{TreeSequenceInitialSegmentVertexSetSubset} and the
vertexwise clause of \(S\)'s \noderef{TreeColoringInvariant}
give equality of the old and new missing sets at every vertex in that
difference. After rewriting \(R.\mathit{coloring}=\sigma\), taking
their finite unions proves \(M_{0,P}=M_{0,R}\).

The same equality \(A_P=A_R\) identifies the numerical bounds
\(|A|\le r<|A|+j\) in the two exact edge-color sets.
\(S\)'s fixed-through-\(j\) field identifies the old and new value
of each individual edge with an index in that range.
Hence, for every color, exactly the same indices witness membership in
the two sets, and \(U_{j,P}=U_{j,R}\).
In particular \(\beta\notin\Gamma_R\) and
\(\beta\notin U_{j,R}\).

If the first input alternative holds, the equality of \(M_0\) and
\(D\) transports its two exclusions to \(R\). Together with the
two exclusions just established, these are exactly the first output
alternative.
If the second input alternative holds, transport its membership in
\(D\) and its Gamma-pair disjointness by the equalities of \(D\),
the pair, and \(U_j\). The beta part of \(S\)'s
missing-two-before field, rewritten by \(R.\mathit{coloring}=\sigma\),
gives the required new missingness on \(B\). These are exactly the
three clauses of the second output alternative.
This is equation (9) for the full original-relative switched
presentation, without changing a parallel edge identity.
\end{proof}
